\documentclass[aip,jcp,reprint,noshowkeys,superscriptaddress]{revtex4-1}
\usepackage{graphicx,dcolumn,bm,xcolor,microtype,multirow,amscd,amsmath,amssymb,amsfonts,physics,longtable,wrapfig}
\usepackage[version=4]{mhchem}

\usepackage[utf8]{inputenc}
\usepackage[T1]{fontenc}
\usepackage{txfonts}

\usepackage[
    colorlinks,
    linkcolor={red!50!black},
    citecolor={red!70!black},
    urlcolor={red!80!black}
	]{hyperref}
\urlstyle{same}

\newcommand{\ie}{\textit{i.e.}}
\newcommand{\eg}{\textit{e.g.}}
\newcommand{\alert}[1]{\textcolor{red}{#1}}
\usepackage[normalem]{ulem}
\newcommand{\titou}[1]{\textcolor{red}{#1}}
\newcommand{\trashPFL}[1]{\textcolor{\red}{\sout{#1}}}
\newcommand{\PFL}[1]{\titou{(\underline{\bf PFL}: #1)}}

\newcommand{\cre}[1]{\Hat{a}_{#1}^{\dag}}
\newcommand{\ani}[1]{\Hat{a}_{#1}}
\newcommand{\bcre}[1]{\Hat{b}_{#1}^{\dag}}
\newcommand{\bani}[1]{\Hat{b}_{#1}}

\newcommand{\mc}{\multicolumn}
\newcommand{\fnm}{\footnotemark}
\newcommand{\fnt}{\footnotetext}
\newcommand{\tabc}[1]{\multicolumn{1}{c}{#1}}
\newcommand{\QP}{\textsc{quantum package}}
\newcommand{\T}[1]{#1^{\intercal}}

% coordinates
\newcommand{\br}{\boldsymbol{r}}
\newcommand{\bt}{\boldsymbol{t}}
\newcommand{\bx}{\boldsymbol{x}}

\newcommand{\bA}{\boldsymbol{A}}
\newcommand{\bB}{\boldsymbol{B}}
\newcommand{\bC}{\boldsymbol{C}}
\newcommand{\bD}{\boldsymbol{D}}
\newcommand{\bH}{\boldsymbol{H}}
\newcommand{\bI}{\boldsymbol{1}}
\newcommand{\bK}{\boldsymbol{K}}
\newcommand{\bL}{\boldsymbol{L}}
\newcommand{\bM}{\boldsymbol{M}}
\newcommand{\bO}{\boldsymbol{0}}
\newcommand{\bP}{\boldsymbol{P}}
\newcommand{\bQ}{\boldsymbol{Q}}
\newcommand{\bR}{\boldsymbol{R}}
\newcommand{\bS}{\boldsymbol{S}}
\newcommand{\bT}{\boldsymbol{T}}
\newcommand{\bX}{\boldsymbol{X}}
\newcommand{\bY}{\boldsymbol{Y}}
\newcommand{\bZ}{\boldsymbol{Z}}

\newcommand{\hF}{\Hat{F}}
\newcommand{\hH}{\Hat{H}}
\newcommand{\hP}{\Hat{P}}
\newcommand{\hQ}{\Hat{Q}}
\newcommand{\hT}{\Hat{T}}
\newcommand{\hZ}{\Hat{Z}}
\newcommand{\timeorder}{\mathcal{T}}


\newcommand{\bDel}{\boldsymbol{\Delta}}
\newcommand{\bOm}{\boldsymbol{\Omega}}

\newcommand{\cK}{\mathcal{K}}
\newcommand{\cQ}{\mathcal{Q}}

% orbital energies
\newcommand{\eps}{\epsilon}
\newcommand{\irr}{\text{irr}}

% shortcuts for greek letters
\newcommand{\si}{\sigma}
\newcommand{\la}{\lambda}
\newcommand{\ii}{\mathrm{i}}
\newcommand{\up}{\uparrow}
\newcommand{\dw}{\downarrow}
\newcommand{\upup}{\up\up}
\newcommand{\updw}{\up\dw}
\newcommand{\dwup}{\dw\up}
\newcommand{\dwdw}{\dw\downarrow}

\newcommand{\h}{\text{1h}}
\newcommand{\p}{\text{1p}}
\newcommand{\hp}{\text{1h+1p}}
\newcommand{\hhp}{\text{2h1p}}
\newcommand{\hhpp}{\text{2h2p}}
\newcommand{\pph}{\text{2p1h}}
\newcommand{\ph}{\text{ph}}
\newcommand{\eh}{\text{eh}}
\newcommand{\he}{\text{he}}
\newcommand{\ee}{\text{ee}}
\newcommand{\teh}{{\overline{\text{eh}}}}
\newcommand{\pp}{\text{pp}}
\newcommand{\hh}{\text{hh}}
\newcommand{\xc}{\text{xc}}
\newcommand{\Hx}{\text{Hx}}


% addresses
\newcommand{\LCPQ}{Univ Toulouse, CNRS, Laboratoire de Chimie et Physique Quantiques, Toulouse, France}
\newcommand{\UnHam}{Department of Chemistry, University of Hamburg, 22761 Hamburg, Germany; The Hamburg Centre for Ultrafast Imaging (CUI), Hamburg 22761, Germany}
\newcommand{\UWarsaw}{Institute of Theoretical Physics, Faculty of Physics, University of Warsaw, Warsaw, Poland}

\begin{document}


\title{From Parquet Theory to Coupled-Cluster Doubles}

%\author{Pierre-Fran\c{c}ois \surname{Loos}}
%	\email{loos@irsamc.ups-tlse.fr}
%	\affiliation{\LCPQ}
%
%\author{Antoine \surname{Marie}}
%	\affiliation{\UWarsaw}
%
%\author{Johannes T\"olle}
%        \email{johannes.toelle@uni-hamburg.de}
%        \affiliation{\UnHam}
	
\begin{abstract}
Parquet theory and coupled-cluster theory look very different at first sight.
Parquet theory describes the many-electron problem in terms of a four-point
vertex and separates this vertex into three scattering channels.  Coupled
cluster instead works with a correlated wave function and determines its
excitation amplitudes by projecting a similarity-transformed Hamiltonian.
In these notes, we show how the two pictures can be connected.
We first introduce the parquet equations and their three Bethe--Salpeter
equations (BSEs).  For static kernels, each BSE can be rewritten as a Riccati
equation by introducing $\bT=\bY\cdot\bX^{-1}$.  We then show that this
Riccati equation is itself a decoupling condition after a block-triangular
similarity transformation.  The matrix $\bT$ has a simple physical meaning:
it describes pairs of elementary electron-hole or particle-particle
excitations in the correlated quasiboson vacuum.  Once these pairs are
written in terms of fermionic creation and annihilation operators, they become
doubly excited determinants.  This naturally leads to a single
antisymmetric doubles amplitude $t_{ij}^{ab}$.  The direct electron-hole,
transverse electron-hole, and particle-particle channels then generate the
ring, crossed-ring, and ladder parts of the CCD equations.  When the three
channel kernels are kept equal to the bare Coulomb interaction, the resulting
common-amplitude equation is the CCD doubles equation.  We carefully
distinguish the interaction-generated pieces of this Riccati residual from
the genuine frequency-dependent parquet vertices.  This distinction also
makes clear what additional information is required for a controlled
extension beyond CCD.
\end{abstract}

\maketitle

%%%%%%%%%%%%%%%%%%%%%%%%%
\section{Introduction}
%%%%%%%%%%%%%%%%%%%%%%%%%

Parquet theory and coupled-cluster (CC) theory provide two different ways of
organizing the same many-electron problem.  In parquet theory, the central
object is the full four-point vertex $F$.  Its diagrams are sorted according
to the two-particle channel in which they can be cut.  In CC theory, the
central object is the cluster operator, and the amplitudes are found by
requiring the similarity-transformed Hamiltonian not to couple the reference
determinant to excited determinants.

These two languages are usually developed separately.  Nevertheless, the
structure of the CCD equations already hints at a connection.  CCD contains
ring, crossed-ring, and ladder terms.  Parquet theory contains the direct
electron-hole ($\eh$), transverse electron-hole ($\teh$), and
particle-particle ($\pp$) channels.  The correspondence
\begin{equation}
    \eh \longleftrightarrow \text{ring},
    \qquad
    \teh \longleftrightarrow \text{crossed ring},
    \qquad
    \pp \longleftrightarrow \text{ladder}
    \label{eq:channel_correspondence}
\end{equation}
is therefore very tempting, but the connection is deeper than a simple
comparison of diagrams.

The route followed below is the following.  We start from the parquet
equations and write the three static BSEs.  Each BSE can be transformed into a
Riccati equation.  We then show that this Riccati equation is a decoupling
condition, in much the same way as the CC amplitude equations.  The Riccati
matrix $\bT=\bY\cdot\bX^{-1}$ describes pairs of elementary excitations in
the quasiboson vacuum.  A pair of electron-hole excitations acting on the
reference determinant is simply a doubly excited determinant.  This is the
step that turns a BSE amplitude into a CC doubles amplitude.  Finally, the
three channels are tied together by requiring them to use the same
antisymmetric tensor $t_{ij}^{ab}$.

It is useful to keep one point in mind from the beginning.  The parquet
decomposition and the three channel BSEs are exact.  For a static kernel, the
transformation from the BSE eigenvalue problem to a Riccati equation is also
an exact algebraic step.  The extra ingredient that makes contact with CCD is
the requirement that the three Riccati amplitudes are three views of one and
the same fermionic tensor $t_{ij}^{ab}$.  We will refer to this requirement as
the \emph{common-amplitude closure}.  With bare Coulomb kernels, it gives the
CCD doubles equation.  With dressed kernels, the common-amplitude residual
remains well defined, but it must not be confused with a full parquet update:
the latter requires the genuine four-point channel-reducible vertices.  We
will return to this distinction at the end of the notes.

%%%%%%%%%%%%%%%%%%%%%%%%%
\section{Parquet theory}
\label{sec:parquet}
%%%%%%%%%%%%%%%%%%%%%%%%%

We begin with the usual one- and two-particle quantities.  Numerical labels
denote collective space-spin-time variables,
\begin{equation}
    1\equiv(\bx_1,t_1)=(\br_1,\sigma_1,t_1)
\end{equation}
and repeated labels are integrated.

%%%%%%%%%%%%%%%%%%%%%%%%%
\subsection{One- and two-particle Green's functions}
%%%%%%%%%%%%%%%%%%%%%%%%%

The time-ordered one-particle Green's function is
\begin{equation}
    G(11')
    =
    -\ii
    \mel{\Psi_0^N}{
        \timeorder[
            \hat{\psi}(1)
            \hat{\psi}^{\dagger}(1')
        ]
    }{\Psi_0^N}
    \label{eq:G_def}
\end{equation}
and satisfies the Dyson equation
\begin{equation}
    G(11')
    =
    G_0(11')
    +
    G_0(12)
    \Sigma(22')
    G(2'1').
    \label{eq:Dyson}
\end{equation}
The corresponding two-particle Green's function is
\begin{equation}
    \begin{split}
    G_2(12;1'2')
    =
    (-\ii)^2
    \mel{\Psi_0^N}{
        \timeorder[
            \hat{\psi}(1)
            \hat{\psi}(2)
            \hat{\psi}^{\dagger}(2')
            \hat{\psi}^{\dagger}(1')
        ]
    }{\Psi_0^N}.
    \end{split}
    \label{eq:G2_def}
\end{equation}
Its connected part defines the full four-point vertex $F$.  With the present
convention,
\begin{multline}
    G_2(12;34)
    =
    G(13)G(24)
    -
    G(14)G(23)
    \\
    -
    G(11')G(22')
    F(1'2';3'4')
    G(3'3)G(4'4).
    \label{eq:G2_F}
\end{multline}
The first two terms describe independent propagation of the two fermions.
Everything that connects the two lines is contained in $F$.

Because the particles are fermions, $F$ changes sign when either pair of
external legs is exchanged,
\begin{equation}
    F(12;34)
    =
    -F(21;34)
    =
    -F(12;43)
    =
    F(21;43).
    \label{eq:F_crossing}
\end{equation}
This simple crossing property will later be responsible for the link between
the direct and transverse electron-hole channels.

%%%%%%%%%%%%%%%%%%%%%%%%%
\subsection{The three scattering channels}
%%%%%%%%%%%%%%%%%%%%%%%%%

A connected four-point diagram is called two-particle reducible if it falls
apart after cutting two one-particle propagator lines.  There are three
inequivalent ways of making this cut.  They define the direct electron-hole,
transverse electron-hole, and particle-particle channels.  We denote the
parts of the full vertex reducible in these channels by
$\Phi^\eh$, $\Phi^{\teh}$, and $\Phi^\pp$, respectively.  Everything that is
irreducible in all three channels is collected in the fully irreducible
vertex $\Lambda$.

The full vertex can therefore be written exactly as
\begin{equation}
    F
    =
    \Lambda
    +
    \Phi^\eh
    +
    \Phi^{\teh}
    +
    \Phi^\pp.
    \label{eq:parquet_decomposition}
\end{equation}
From this decomposition, the vertex irreducible in one particular channel is
obtained by removing only the piece reducible in that channel,
\begin{align}
    \Gamma^\eh
    &=
    \Lambda+\Phi^{\teh}+\Phi^\pp,
    \label{eq:Gamma_eh}
    \\
    \Gamma^{\teh}
    &=
    \Lambda+\Phi^\eh+\Phi^\pp,
    \label{eq:Gamma_teh}
    \\
    \Gamma^\pp
    &=
    \Lambda+\Phi^\eh+\Phi^{\teh}.
    \label{eq:Gamma_pp}
\end{align}
Thus, in any channel $r$,
\begin{equation}
    F=\Gamma^r+\Phi^r.
    \label{eq:F_Gamma_Phi}
\end{equation}
This last relation is particularly useful: $\Gamma^r$ is the starting
interaction in channel $r$, while $\Phi^r$ is the extra scattering generated
by iterating that interaction in the same channel.

%%%%%%%%%%%%%%%%%%%%%%%%%
\subsection{Bethe--Salpeter equations}
%%%%%%%%%%%%%%%%%%%%%%%%%

The previous statement is made explicit by the BSEs.  In compact notation,
all three have the same form,
\begin{equation}
    F
    =
    \Gamma^r
    +
    \Gamma^r L_0^r F,
    \label{eq:BSE_symbolic}
\end{equation}
where $L_0^r$ is the independent two-particle propagator in channel $r$.
Equivalently,
\begin{equation}
    \Phi^r
    =
    \Gamma^r L_0^r F.
    \label{eq:Phi_BSE_symbolic}
\end{equation}
For a static kernel, Eq.~\eqref{eq:BSE_symbolic} can be solved algebraically.
Using
\begin{equation}
    F
    =
    \left(1-\Gamma^r L_0^r\right)^{-1}\Gamma^r,
\end{equation}
Eq.~\eqref{eq:Phi_BSE_symbolic} becomes
\begin{equation}
    \Phi^r(\omega)
    =
    \Gamma^r
    \left[
        \left(L_0^r\right)^{-1}(\omega)-\Gamma^r
    \right]^{-1}
    \Gamma^r.
    \label{eq:Phi_resolvent_exact}
\end{equation}
Equation~\eqref{eq:Phi_resolvent_exact} is the genuine channel-reducible
four-point vertex.  Even when $\Gamma^r$ is static, $\Phi^r$ remains a
frequency-dependent object through the BSE resolvent.  This distinction will
be important below: the interaction-generated part of a Riccati residual is a
reduced amplitude-space quantity and is not identical to
$\Phi^r(\omega)$.

The three channels differ only by the way the external legs are paired and by
the corresponding signs and symmetry factors.

For reference, the direct electron-hole equation can be written as
\begin{multline}
    F(12;34)
    =
    \Gamma^\eh(12;34)
    \\
    +
    \Gamma^\eh(13';31')
    G(1'4')G(2'3')
    F(4'2;2'4),
    \label{eq:BSE_eh}
\end{multline}
whereas the transverse electron-hole channel reads
\begin{multline}
    F(12;34)
    =
    \Gamma^{\teh}(12;34)
    \\
    -
    \Gamma^{\teh}(3'2;32')
    G(1'3')G(2'4')
    F(14';1'4),
    \label{eq:BSE_teh}
\end{multline}
and the particle-particle channel is
\begin{multline}
    F(12;34)
    =
    \Gamma^\pp(12;34)
    \\
    -
    \frac{1}{2}
    \Gamma^\pp(12;1'2')
    G(1'3')G(2'4')
    F(3'4';34).
    \label{eq:BSE_pp}
\end{multline}
The minus sign in the transverse channel comes from crossing two fermionic
legs.  The factor $1/2$ in the particle-particle channel avoids counting an
antisymmetric pair twice.

The parquet equations are obtained by solving these three BSEs together with
Eqs.~\eqref{eq:Gamma_eh}--\eqref{eq:Gamma_pp}.  To close the equations, one
must still choose the fully irreducible vertex.  In the parquet approximation,
we take
\begin{equation}
    \Lambda\approx\bar v,
    \label{eq:parquet_approximation}
\end{equation}
where $\bar v$ is the antisymmetrized Coulomb interaction.  The simplest
starting point is even more restrictive: all three channel kernels are kept
equal to the bare interaction,
\begin{equation}
    \Gamma^\eh
    =
    \Gamma^{\teh}
    =
    \Gamma^\pp
    =
    \Lambda
    =
    \bar v.
    \label{eq:bare_kernel_limit}
\end{equation}
We will see that this is precisely the limit in which the connection with CCD
becomes exact.

%%%%%%%%%%%%%%%%%%%%%%%%%
\section{Static Bethe--Salpeter equations}
\label{sec:static_bse}
%%%%%%%%%%%%%%%%%%%%%%%%%

We now move to an orbital representation and assume that the channel kernels
are static.  Occupied spin orbitals are denoted by $i,j,k,l$, virtual spin
orbitals by $a,b,c,d$, and general spin orbitals by $p,q,r,s$.  The reference
one-particle propagator is taken diagonal, with orbital energies $\epsilon_p$.

The free two-particle propagator already contains the energy differences that
will later appear in the CCD equation.  For example, the forward direct
$\eh$ propagator is
\begin{equation}
    L_{0,ia}^{\eh,+}(\omega)
    =
    \frac{1}{
        \omega-(\epsilon_a-\epsilon_i)+\ii\eta
    }.
    \label{eq:L0_eh_forward}
\end{equation}
Its inverse therefore contains $\epsilon_a-\epsilon_i$.  This point will be
important when we identify the origin of the usual doubles energy
denominator.

%%%%%%%%%%%%%%%%%%%%%%%%%
\subsection{Direct electron-hole channel}
%%%%%%%%%%%%%%%%%%%%%%%%%

For a static kernel, the direct $\eh$ BSE becomes
\begin{equation}
    \begin{pmatrix}
        \bA^\eh & \bB^\eh
        \\
        -{\bB^\eh}^{\dagger} & -{\bA^\eh}^{\dagger}
    \end{pmatrix}
    \cdot
    \begin{pmatrix}
        \bX^\eh
        \\
        \bY^\eh
    \end{pmatrix}
    =
    \begin{pmatrix}
        \bX^\eh
        \\
        \bY^\eh
    \end{pmatrix}
    \cdot
    \bOm^\eh.
    \label{eq:BSE_matrix_eh}
\end{equation}
We split
\begin{equation}
    \bA^\eh
    =
    \bA_0^\eh+\bK_A^\eh,
    \label{eq:A_split_eh}
\end{equation}
with
\begin{equation}
    [\bA_0^\eh]_{ia,jb}
    =
    (\epsilon_a-\epsilon_i)
    \delta_{ij}\delta_{ab}.
    \label{eq:A0_eh}
\end{equation}
For the bare antisymmetrized Coulomb kernel,
\begin{align}
    [\bA^\eh]_{ia,jb}
    &=
    (\epsilon_a-\epsilon_i)
    \delta_{ij}\delta_{ab}
    +
    \mel{aj}{}{ib},
    \label{eq:A_eh_bare}
    \\
    [\bB^\eh]_{ia,jb}
    &=
    \mel{ab}{}{ij}.
    \label{eq:B_eh_bare}
\end{align}

%%%%%%%%%%%%%%%%%%%%%%%%%
\subsection{Transverse electron-hole channel}
%%%%%%%%%%%%%%%%%%%%%%%%%

The transverse channel pairs the external legs in the crossed order.  Using
the composite indices $(ib,ja)$, we write
\begin{equation}
    \begin{pmatrix}
        \bA^{\teh} & -\bB^{\teh}
        \\
        {\bB^{\teh}}^{\dagger} & -{\bA^{\teh}}^{\dagger}
    \end{pmatrix}
    \cdot
    \begin{pmatrix}
        \bX^{\teh}
        \\
        \bY^{\teh}
    \end{pmatrix}
    =
    \begin{pmatrix}
        \bX^{\teh}
        \\
        \bY^{\teh}
    \end{pmatrix}
    \cdot
    \bOm^{\teh}.
    \label{eq:BSE_matrix_teh}
\end{equation}
The sign changes compared with Eq.~\eqref{eq:BSE_matrix_eh} come directly from
crossing.  They will later generate the crossed-ring terms of CCD.  To make
the sign bookkeeping explicit, we split
\begin{equation}
    \bA^{\teh}
    =
    \bA_0^{\teh}+\bK_A^{\teh},
    \label{eq:A_split_teh}
\end{equation}
with, in the crossed composite ordering,
\begin{equation}
    [\bA_0^{\teh}]_{ib,kc}
    =
    (\epsilon_b-\epsilon_i)\delta_{ik}\delta_{bc}.
    \label{eq:A0_teh}
\end{equation}
For the bare antisymmetrized Coulomb kernel we use
\begin{align}
    [\bK_A^{\teh}]_{ib,kc}
    &=
    \mel{bk}{}{ic},
    \label{eq:KA_teh_bare}
    \\
    [\bB^{\teh}]_{ib,ja}
    &=
    \mel{ba}{}{ij}
    =
    -\mel{ab}{}{ij}.
    \label{eq:B_teh_bare}
\end{align}
Thus the crossing signs are carried by the interaction blocks themselves;
they do not reverse the sign of the free propagation.

%%%%%%%%%%%%%%%%%%%%%%%%%
\subsection{Particle-particle channel}
%%%%%%%%%%%%%%%%%%%%%%%%%

In the $\pp$ channel, the forward sector contains two particles and the
backward sector contains two holes.  We write
\begin{equation}
    \begin{pmatrix}
        \bC^\pp & \bB^\pp
        \\
        -{\bB^\pp}^{\dagger} & -{\bD^\pp}^{\dagger}
    \end{pmatrix}
    \cdot
    \begin{pmatrix}
        \bX^\pp
        \\
        \bY^\pp
    \end{pmatrix}
    =
    \begin{pmatrix}
        \bX^\pp
        \\
        \bY^\pp
    \end{pmatrix}
    \cdot
    \bOm^\pp.
    \label{eq:BSE_matrix_pp}
\end{equation}
The free parts of $\bC^\pp$ and $\bD^\pp$ contain the two-particle and
two-hole energies.  Antisymmetric pair normalization is understood throughout.

%%%%%%%%%%%%%%%%%%%%%%%%%
\section{From a BSE to an amplitude equation}
\label{sec:riccati}
%%%%%%%%%%%%%%%%%%%%%%%%%

The first real step toward coupled cluster is purely algebraic.  Consider the
direct $\eh$ problem.  Its two block equations are
\begin{align}
    \bA^\eh\cdot\bX^\eh
    +
    \bB^\eh\cdot\bY^\eh
    &=
    \bX^\eh\cdot\bOm^\eh,
    \label{eq:BSE_upper}
    \\
    -{\bB^\eh}^{\dagger}\cdot\bX^\eh
    -
    {\bA^\eh}^{\dagger}\cdot\bY^\eh
    &=
    \bY^\eh\cdot\bOm^\eh.
    \label{eq:BSE_lower}
\end{align}
Assuming that $\bX^\eh$ is invertible in the relevant subspace, define
\begin{equation}
    \bT^\eh
    =
    \bY^\eh\cdot(\bX^\eh)^{-1}.
    \label{eq:T_eh_YX}
\end{equation}
The upper equation gives
\begin{equation}
    \bX^\eh\cdot\bOm^\eh\cdot(\bX^\eh)^{-1}
    =
    \bA^\eh+\bB^\eh\cdot\bT^\eh.
    \label{eq:Omega_similarity}
\end{equation}
Inserting this result into the lower equation removes both $\bX^\eh$ and
$\bOm^\eh$ and gives
\begin{equation}
    {\bB^\eh}^{\dagger}
    +
    {\bA^\eh}^{\dagger}\cdot\bT^\eh
    +
    \bT^\eh\cdot\bA^\eh
    +
    \bT^\eh\cdot\bB^\eh\cdot\bT^\eh
    =
    \bO.
    \label{eq:Riccati_eh}
\end{equation}
The BSE eigenvalue problem has therefore become a nonlinear equation for the
matrix $\bT^\eh$.

Exactly the same manipulation can be made in the other two channels.  For the
transverse $\eh$ problem,
\begin{equation}
    \bT^{\teh}
    =
    \bY^{\teh}\cdot(\bX^{\teh})^{-1}
\end{equation}
gives
\begin{equation}
    -{\bB^{\teh}}^{\dagger}
    +
    {\bA^{\teh}}^{\dagger}\cdot\bT^{\teh}
    +
    \bT^{\teh}\cdot\bA^{\teh}
    -
    \bT^{\teh}\cdot\bB^{\teh}\cdot\bT^{\teh}
    =
    \bO.
    \label{eq:Riccati_teh}
\end{equation}
For the $\pp$ problem, the natural Riccati matrix maps the forward vv space
onto the backward oo space,
\begin{equation}
    \bT^\pp
    =
    \bY^\pp\cdot(\bX^\pp)^{-1},
    \label{eq:T_pp_YX}
\end{equation}
and obeys
\begin{equation}
    {\bB^\pp}^{\dagger}
    +
    {\bD^\pp}^{\dagger}\cdot\bT^\pp
    +
    \bT^\pp\cdot\bC^\pp
    +
    \bT^\pp\cdot\bB^\pp\cdot\bT^\pp
    =
    \bO.
    \label{eq:Riccati_pp}
\end{equation}
We have thus replaced the three BSE eigenvalue problems by three amplitude
equations.

%%%%%%%%%%%%%%%%%%%%%%%%%
\subsection{The Riccati equation is a decoupling condition}
%%%%%%%%%%%%%%%%%%%%%%%%%

The previous algebra becomes more meaningful if we look at it as a similarity
transformation.  Introduce
\begin{equation}
    \bS(\bT)
    =
    \begin{pmatrix}
        \bI & \bO
        \\
        \bT & \bI
    \end{pmatrix},
    \qquad
    \bS^{-1}(\bT)
    =
    \begin{pmatrix}
        \bI & \bO
        \\
        -\bT & \bI
    \end{pmatrix}.
    \label{eq:S_T}
\end{equation}
For the direct $\eh$ BSE matrix
\begin{equation}
    \bH^\eh
    =
    \begin{pmatrix}
        \bA^\eh & \bB^\eh
        \\
        -{\bB^\eh}^{\dagger} & -{\bA^\eh}^{\dagger}
    \end{pmatrix},
\end{equation}
one finds
\begin{equation}
    \bS^{-1}\cdot\bH^\eh\cdot\bS
    =
    \begin{pmatrix}
        \bA^\eh+\bB^\eh\cdot\bT
        &
        \bB^\eh
        \\
        -\bR^\eh[\bT]
        &
        -{\bA^\eh}^{\dagger}-\bT\cdot\bB^\eh
    \end{pmatrix},
    \label{eq:block_similarity}
\end{equation}
where
\begin{equation}
    \bR^\eh[\bT]
    =
    {\bB^\eh}^{\dagger}
    +
    {\bA^\eh}^{\dagger}\cdot\bT
    +
    \bT\cdot\bA^\eh
    +
    \bT\cdot\bB^\eh\cdot\bT.
    \label{eq:Riccati_residual}
\end{equation}
The Riccati equation simply says that the lower-left block of the transformed
matrix must vanish.

If $\bP_r$ and $\bQ_r$ project onto the forward and backward sectors of
channel $r$, this condition can be written as
\begin{equation}
    \bQ_r
    \cdot
    \bS_r^{-1}(\bT^r)
    \cdot
    \bH^r
    \cdot
    \bS_r(\bT^r)
    \cdot
    \bP_r
    =
    \bO.
    \label{eq:BSE_decoupling}
\end{equation}
This is our first direct contact with coupled cluster.  The CCD amplitude
equation can also be written as a decoupling condition,
\begin{equation}
    \hQ_2
    e^{-\hT_2}
    \hH
    e^{\hT_2}
    \hP
    =
    0,
    \label{eq:CCD_decoupling}
\end{equation}
where $\hP=\dyad{\Phi}{\Phi}$ projects onto the reference determinant and
$\hQ_2$ projects onto the doubly excited determinants.  The two equations do
not yet live in the same spaces, but they have exactly the same purpose: a
similarity transformation is chosen so that a coupling between two sectors
vanishes.

%%%%%%%%%%%%%%%%%%%%%%%%%
\section{What does \texorpdfstring{$\bT=\bY\cdot\bX^{-1}$}{T = Y times X inverse} describe?}
\label{sec:boson_to_doubles}
%%%%%%%%%%%%%%%%%%%%%%%%%

The remaining step becomes clear in operator language.  Define an elementary
electron-hole excitation operator
\begin{equation}
    \bcre{ia}
    =
    \cre{a}\ani{i},
    \qquad
    \bani{ia}
    =
    \cre{i}\ani{a}.
    \label{eq:B_ph_creation}
\end{equation}
Within the quasiboson approximation, these composite fermionic operators are
treated as bosonic creation and annihilation operators.  An RPA excitation
operator then has the form
\begin{equation}
    \hQ_n^{\dagger}
    =
    \sum_{ia}
    \left(
        X_{ia,n}\bcre{ia}
        -
        Y_{ia,n}\bani{ia}
    \right).
    \label{eq:RPA_excitation_operator}
\end{equation}
Because $\bY$ is nonzero, the RPA vacuum is not the uncorrelated reference.  In
the quasiboson picture it can be written as
\begin{equation}
    \ket{\Psi_{\rm RPA}}
    =
    e^{\hZ}\ket{\Phi},
    \label{eq:RPA_vacuum}
\end{equation}
with
\begin{equation}
    \hZ
    =
    \frac{1}{2}
    \sum_{iajb}
    Z_{ia,jb}
    \bcre{ia}\bcre{jb}.
    \label{eq:Z_operator}
\end{equation}
The condition that every RPA annihilation operator kills this vacuum,
\begin{equation}
    \hQ_n\ket{\Psi_{\rm RPA}}=0,
\end{equation}
gives
\begin{equation}
    \bZ
    =
    \bY\cdot\bX^{-1}
    =
    \bT.
    \label{eq:Z_equals_T}
\end{equation}
The Riccati matrix $\bT$ is therefore the amplitude for creating a \emph{pair}
of elementary quasibosonic excitations in the correlated vacuum.

Now return to the exact fermionic operators.  Two electron-hole excitations
acting on the reference give
\begin{equation}
    \begin{split}
    \bcre{ia}\bcre{jb}\ket{\Phi}
    &=
    \cre{a}\ani{i}\cre{b}\ani{j}\ket{\Phi}
    \\
    &=
    \cre{a}\cre{b}\ani{j}\ani{i}\ket{\Phi}
    \\
    &=
    \ket{\Phi_{ij}^{ab}},
    \end{split}
    \label{eq:ph_pair_to_double}
\end{equation}
up to the fixed operator-ordering convention.  In other words, a pair of
$\eh$ excitations is exactly a $\hhpp$ excitation.  The pair index
$(ia,jb)$ can therefore be unfolded into the four fermionic indices $(ijab)$.
This is the missing link between the BSE and CCD spaces.

The same idea can be expressed directly at the level of the decoupling
condition.  In the quasiboson picture, the Riccati equation is equivalent to
requiring the pair-creation part of the transformed RPA Hamiltonian to vanish,
\begin{equation}
    \mel{\Phi}{
        \bani{jb}\bani{ia}
        e^{-\hZ}
        \hH_{\rm RPA}
        e^{\hZ}
    }{\Phi}
    =
    0.
    \label{eq:RPA_pair_decoupling}
\end{equation}
After the pair $\bcre{ia}\bcre{jb}$ is interpreted as a true fermionic double
excitation, this has the same form as the CCD projection
\begin{equation}
    \mel{\Phi_{ij}^{ab}}{
        e^{-\hT_2}
        \hH
        e^{\hT_2}
    }{\Phi}
    =
    0.
    \label{eq:CCD_projection}
\end{equation}
The difference is that the first equation was obtained in a quasiboson space,
whereas the second one acts in the fermionic Hilbert space.  The rest of the
construction is precisely about restoring that fermionic structure.

%%%%%%%%%%%%%%%%%%%%%%%%%
\section{One doubles amplitude for the three channels}
\label{sec:common_t}
%%%%%%%%%%%%%%%%%%%%%%%%%

If the three BSEs were solved independently, each channel would have its own
Riccati matrix and its own quasiboson vacuum.  CCD does not work this way.
There is only one doubles operator,
\begin{equation}
    \hT_2
    =
    \frac{1}{4}
    \sum_{ijab}
    t_{ij}^{ab}
    \cre{a}\cre{b}\ani{j}\ani{i},
    \label{eq:T2_def}
\end{equation}
with
\begin{equation}
    t_{ij}^{ab}
    =
    -t_{ji}^{ab}
    =
    -t_{ij}^{ba}
    =
    t_{ji}^{ba}.
    \label{eq:t_antisymmetry}
\end{equation}
The key step is therefore to use this same tensor in the three channel
representations,
\begin{align}
    T^\eh_{ia,jb}
    &=
    t_{ij}^{ab},
    \label{eq:T_eh_map}
    \\
    T^{\teh}_{ib,ja}
    &=
    t_{ij}^{ab},
    \label{eq:T_teh_map}
    \\
    T^\pp_{ij,ab}
    &=
    t_{ij}^{ab}.
    \label{eq:T_pp_map}
\end{align}
The last convention follows the natural orientation
$\bT^\pp = \bY^\pp \cdot (\bX^\pp)^{-1}$, which maps the forward vv space onto the
backward oo space.

These three matrices are not three amplitudes.  They are three ways of
arranging the indices of the same fermionic tensor.  This common-amplitude
condition is the additional information that is absent when the three RPA
problems are solved separately.

The transverse channel has a particularly simple role in this picture.  A
single direct $\eh$ quasiboson amplitude is symmetric under exchange of its
composite bosonic indices.  This is not enough to guarantee separate
antisymmetry in the occupied and virtual fermionic indices.  The crossed
pairing supplies the missing exchange.  Schematically,
\begin{equation}
    t_{ij}^{ab}
    \sim
    T^\eh_{ia,jb}
    -
    T^\eh_{ib,ja}.
    \label{eq:crossed_antisymmetry}
\end{equation}
Thus the transverse channel is not an optional extra term.  It is what allows
the ring-like quasiboson amplitudes to be combined into an antisymmetric
fermionic doubles tensor.  The $\pp$ representation describes ladder
propagation of that same tensor.

%%%%%%%%%%%%%%%%%%%%%%%%%
\section{From the three parquet channels to the common amplitude equation}
\label{sec:amplitude_equation}
%%%%%%%%%%%%%%%%%%%%%%%%%

We can now put the pieces together.  The derivation is easiest to follow if we
first separate, in every channel Riccati equation, the free propagation from
the interaction terms.

%%%%%%%%%%%%%%%%%%%%%%%%%
\subsection{The free term comes from \texorpdfstring{$L_0^{-1}$}{the inverse free propagator}}
%%%%%%%%%%%%%%%%%%%%%%%%%

Consider again the direct $\eh$ equation and use
$\bA^\eh=\bA_0^\eh+\bK_A^\eh$.  The part containing only the free propagator is
\begin{equation}
    \bA_0^\eh\cdot\bT^\eh
    +
    \bT^\eh\cdot\bA_0^\eh.
    \label{eq:Delta_matrix}
\end{equation}
With $T^\eh_{ia,jb}=t_{ij}^{ab}$, this becomes
\begin{equation}
    \Delta_{ij}^{ab}[\bt]
    =
    (\epsilon_a+\epsilon_b-\epsilon_i-\epsilon_j)
    t_{ij}^{ab}.
    \label{eq:Delta_def}
\end{equation}
This term is sometimes introduced from the CC side as the orbital-energy part
of the doubles equation.  Here its origin is more direct: it is simply the
inverse free two-particle propagator acting on the pair amplitude,
\begin{equation}
    \bDel[\bt]
    \longleftrightarrow
    (\bL_0^r)^{-1}\cdot\bT^r.
    \label{eq:Delta_L0}
\end{equation}
Once written in CC language, the same contribution comes from the one-body
commutator $\comm{\hF}{\hT_2}$.

The $\pp$ channel gives the same total energy difference by combining the
particle-pair and hole-pair energies.  The transverse $\eh$ channel does so
as well.  Indeed, using Eqs.~\eqref{eq:A0_teh} and
\eqref{eq:T_teh_map},
\begin{align}
    &\left[
        {\bA_0^{\teh}}^{\dagger}\cdot\bT^{\teh}
        +
        \bT^{\teh}\cdot\bA_0^{\teh}
    \right]_{ib,ja}
    \nonumber\\
    &\qquad=
    \left[
        (\epsilon_b-\epsilon_i)
        +
        (\epsilon_a-\epsilon_j)
    \right]
    t_{ij}^{ab}
    =
    \Delta_{ij}^{ab}[\bt].
    \label{eq:Delta_teh}
\end{align}
The crossing signs in Eq.~\eqref{eq:Riccati_teh} therefore reside in the
kernel blocks $- {\bB^{\teh}}^{\dagger}$ and
$-\bT^{\teh}\bB^{\teh}\bT^{\teh}$, together with the crossed orbital
ordering of $\bK_A^{\teh}$; they do not multiply the free term by an
overall minus sign.  After unfolding to $ijab$ notation, all three channels
contain the same $+\Delta[\bt]$ contribution.

%%%%%%%%%%%%%%%%%%%%%%%%%
\subsection{What each Riccati equation adds}
%%%%%%%%%%%%%%%%%%%%%%%%%

For fixed kernels, let $\cQ_r[\bt;\Gamma^r]$
denote the Riccati residual of channel $r$ after its composite indices have
been unfolded into $ijab$ notation.  At $\bt = \bO$, only the block of the
kernel that directly couples the forward and backward sectors remains.  We
denote this channel driver by
\begin{equation}
    \cK^r
    =
    \cQ_r[\bO;\Gamma^r].
    \label{eq:K_r_def}
\end{equation}
Each residual also contains the same free contribution $\Delta[\bt]$.
We therefore define the interaction-generated Riccati contribution
\begin{equation}
    \Xi^r[\bt;\Gamma^r]
    =
    \cQ_r[\bt;\Gamma^r]
    -
    \cK^r
    -
    \Delta[\bt].
    \label{eq:Xi_r_def}
\end{equation}
Equivalently,
\begin{equation}
    \cQ_r[\bt;\Gamma^r]
    =
    \cK^r
    +
    \Delta[\bt]
    +
    \Xi^r[\bt;\Gamma^r].
    \label{eq:Q_r_decomposition}
\end{equation}
The notation $\Xi^r$ is deliberate.  The genuine parquet quantity
$\Phi^r(\omega)$ is the four-point vertex of
Eq.~\eqref{eq:Phi_resolvent_exact}; by contrast,
$\Xi^r[\bt;\Gamma^r]$ is the interaction part of the Riccati decoupling
condition after reduction to the common amplitude space.  The two objects are
related through the same channel BSE, but they are not identical and should
not be denoted by the same symbol.

At the bare-kernel level, the three $\Xi^r$ take a simple form.  The direct
$\eh$ channel
gives
\begin{equation}
    \begin{split}
    [\Xi^\eh]_{ij}^{ab}[\bt]
    ={}
    &
    \sum_{kc}
    \mel{ak}{}{ic}
    t_{kj}^{cb}
    +
    \sum_{kc}
    \mel{bk}{}{jc}
    t_{ki}^{ca}
    \\
    &+
    \sum_{klcd}
    t_{ki}^{ca}
    \mel{kl}{}{cd}
    t_{lj}^{db},
    \end{split}
    \label{eq:Xi_eh}
\end{equation}
which is the ring part.  The transverse channel gives
\begin{equation}
    \begin{split}
    [\Xi^{\teh}]_{ij}^{ab}[\bt]
    ={}
    &-
    \sum_{kc}
    \mel{kb}{}{ic}
    t_{kj}^{ac}
    -
    \sum_{kc}
    \mel{ka}{}{jc}
    t_{ki}^{bc}
    \\
    &-
    \sum_{klcd}
    t_{ki}^{bc}
    \mel{kl}{}{cd}
    t_{lj}^{ad},
    \end{split}
    \label{eq:Xi_teh}
\end{equation}
which is the crossed-ring part.  Finally, the $\pp$ channel gives
\begin{equation}
    \begin{split}
    [\Xi^\pp]_{ij}^{ab}[\bt]
    ={}
    &
    \frac{1}{2}
    \sum_{cd}
    \mel{ab}{}{cd}
    t_{ij}^{cd}
    +
    \frac{1}{2}
    \sum_{kl}
    \mel{kl}{}{ij}
    t_{kl}^{ab}
    \\
    &+
    \frac{1}{4}
    \sum_{klcd}
    t_{ij}^{cd}
    \mel{cd}{}{kl}
    t_{kl}^{ab},
    \end{split}
    \label{eq:Xi_pp}
\end{equation}
which is the ladder part.  The first two expressions are related by crossing;
the different signs simply reflect the chosen orientation of the external
legs.

%%%%%%%%%%%%%%%%%%%%%%%%%
\subsection{Assembling the common-amplitude residual}
%%%%%%%%%%%%%%%%%%%%%%%%%

The three Riccati equations are three representations of the same fermionic
doubles amplitude, not three residuals that should be added in full.  In
particular, adding them directly would count the same free propagation
$\Delta[\bt]$ three times and would also repeat the channel drivers.  The
common-amplitude closure is instead assembled by including each common piece
once and each channel-specific interaction contribution once.

The fully irreducible interaction supplies the common doubles driver.  The
free propagation supplies a single $\Delta[\bt]$.  The channel-specific parts
are the three $\Xi^r$ defined in Eq.~\eqref{eq:Xi_r_def}.  We therefore define
\begin{equation}
    \mathcal R[\bt]
    =
    \Delta[\bt]
    +
    \Lambda
    +
    \Xi^\eh[\bt;\Gamma^\eh]
    +
    \Xi^{\teh}[\bt;\Gamma^{\teh}]
    +
    \Xi^\pp[\bt;\Gamma^\pp].
    \label{eq:R_common_final}
\end{equation}
The common-amplitude equation is
\begin{equation}
    \mathcal R[\bt]=0.
    \label{eq:common_amplitude_equation}
\end{equation}
Equation~\eqref{eq:R_common_final} is an additional fermionic closure imposed
on the three channel Riccati representations.  It should not be confused with
summing the three BSEs.  Its ingredients have distinct origins: $\Delta$ is
the common free propagation, $\Lambda$ is the common fully irreducible
driver, and the three $\Xi^r$ are the channel-specific interaction terms.
This construction avoids any cancellation between different copies of the
free term.

%%%%%%%%%%%%%%%%%%%%%%%%%
\section{The bare-kernel common-amplitude limit is CCD}
\label{sec:ccd_limit}
%%%%%%%%%%%%%%%%%%%%%%%%%

We now set
\begin{equation}
    \Gamma^\eh
    =
    \Gamma^{\teh}
    =
    \Gamma^\pp
    =
    \Lambda
    =
    \bar v.
\end{equation}
The three Riccati contributions in Eq.~\eqref{eq:R_common_final} are then
exactly the ring, crossed-ring, and ladder contractions written in
Eqs.~\eqref{eq:Xi_eh}--\eqref{eq:Xi_pp}.  The remaining CCD terms are the
one-body mosaic contributions.  We denote them by $\mathcal M[\bt]$.  They
may either be kept explicitly or absorbed into an amplitude-dependent
one-body matrix, which simply replaces the bare free propagation by its
dressed counterpart.

The complete equation is therefore
\begin{equation}
    \begin{split}
    0
    ={}
    &
    \Delta_{ij}^{ab}[\bt]
    +
    \mel{ab}{}{ij}
    +
    \mathcal M_{ij}^{ab}[\bt]
    \\
    &+
    [\Xi^\eh]_{ij}^{ab}[\bt]
    +
    [\Xi^{\teh}]_{ij}^{ab}[\bt]
    +
    [\Xi^\pp]_{ij}^{ab}[\bt].
    \end{split}
    \label{eq:CCD_channel_equation}
\end{equation}
This is the CCD doubles equation written channel by channel.  In the usual CC
language, the same equation is
\begin{equation}
    \mel{\Phi_{ij}^{ab}}{
        e^{-\hT_2}
        \hH
        e^{\hT_2}
    }{\Phi}
    =
    0.
    \label{eq:CCD_exact_projection}
\end{equation}

The connection can now be stated very simply.  The direct $\eh$ BSE provides
the ring terms, the transverse $\eh$ BSE provides the crossed-ring terms, and
the $\pp$ BSE provides the ladder terms.  The three are not solved with three
independent amplitudes.  They are evaluated with the same antisymmetric
$t_{ij}^{ab}$.  The mosaic terms take care of the one-body renormalization.
With these ingredients, the bare-kernel common-amplitude construction
reproduces the CCD doubles equation.

It is worth stressing what would go wrong if the three RPA problems were
solved independently.  In that case one would obtain three different
matrices $\bT^r=\bY^r \cdot (\bX^r)^{-1}$ and therefore three different quasiboson
vacua.  There would be no reason for those three matrices to be different
arrangements of one antisymmetric fermionic tensor.  The equality with CCD
comes from the common-amplitude closure, not from solving three unrelated RPA
problems.

%%%%%%%%%%%%%%%%%%%%%%%%%
\section{Beyond CCD: what a parquet update requires}
\label{sec:beyond_ccd}
%%%%%%%%%%%%%%%%%%%%%%%%%

The distinction between $\Phi^r$ and $\Xi^r$ becomes essential beyond the
bare-kernel construction.  The quantities $\Xi^r[\bt;\Gamma^r]$ are pieces of
the common-amplitude Riccati residual.  They cannot simply be inserted into
the parquet kernel identities
\begin{align}
    \Gamma^\eh
    &=
    \Lambda+\Phi^{\teh}+\Phi^\pp,
    \\
    \Gamma^{\teh}
    &=
    \Lambda+\Phi^\eh+\Phi^\pp,
    \\
    \Gamma^\pp
    &=
    \Lambda+\Phi^\eh+\Phi^{\teh},
\end{align}
because those equations require the genuine four-point reducible vertices.

For a fixed static channel kernel, the required object is obtained from the
channel BSE itself,
\begin{equation}
    \Phi^r(\omega)
    =
    \Gamma^r
    \left[
        \left(L_0^r\right)^{-1}(\omega)-\Gamma^r
    \right]^{-1}
    \Gamma^r,
    \label{eq:Phi_for_kernel_update}
\end{equation}
which is Eq.~\eqref{eq:Phi_resolvent_exact}.  A genuine parquet iteration
would use these $\Phi^r$ objects in the three kernel identities above.
However, Eq.~\eqref{eq:Phi_for_kernel_update} also makes clear that a static
input kernel generally produces a frequency-dependent reducible vertex, so a
full parquet update does not remain within the static-kernel manifold.  A
static implementation would therefore require an additional, explicitly
defined frequency projection or static approximation.

One may instead consider feeding the Riccati quantities $\Xi^r[\bt]$ back
into effective static kernels.  Such a procedure could be useful, but it
would define a distinct CC-parametrized channel-feedback approximation rather
than the standard parquet equations.  Its diagrammatic content and possible
double counting would have to be established separately.  We therefore do
not identify $\Xi^r$ with $\Phi^r$ in the present notes.

%%%%%%%%%%%%%%%%%%%%%%%%%
\section{One-shot quasiparticle energies}
\label{sec:self_energy}
%%%%%%%%%%%%%%%%%%%%%%%%%

If the goal is a quasiparticle spectrum rather than a ground-state correlation
energy, the one-body self-energy can be evaluated after the entire two-body
problem has been converged for a fixed reference propagator.  The natural
order is therefore
\begin{equation}
    G_0
    \longrightarrow
    \text{converged two-body problem}
    \longrightarrow
    \Sigma_{pp}(\omega)
    \longrightarrow
    \epsilon_p^{\rm QP}.
    \label{eq:one_shot_scheme}
\end{equation}
For a one-shot calculation, only the diagonal self-energy elements are needed.
The screened couplings entering $\Sigma_{pp}(\omega)$ can be constructed from
the channel BSE eigenvectors and eigenvalues, or equivalently from the genuine
channel resolvents entering $\Phi^r(\omega)$ in
Eq.~\eqref{eq:Phi_resolvent_exact}.  In the present construction, the two-body
problem is solved first; the spectral information required for the self-energy
is reconstructed only afterwards.

This last step is logically separate from the derivation of the common
amplitude equation.  The main point of the present notes is the two-body
connection between parquet and CCD.  A detailed derivation of the screened
couplings and of the resulting self-energy can therefore be added once the
static two-body construction has been fully settled.

%%%%%%%%%%%%%%%%%%%%%%%%%
\section{Summary}
\label{sec:summary}
%%%%%%%%%%%%%%%%%%%%%%%%%

The connection developed in these notes can be followed as one continuous
sequence,
\begin{equation}
    \begin{split}
    \text{parquet BSE}
    &\longrightarrow
    \text{static BSE matrix}
    \longrightarrow
    \text{Riccati equation}
    \\
    &\longrightarrow
    \text{pair amplitude}
    \longrightarrow
    t_{ij}^{ab}.
    \end{split}
    \label{eq:central_sequence}
\end{equation}
The BSE-to-Riccati step is a similarity-transformation decoupling.  The matrix
$\bT=\bY \cdot \bX^{-1}$ describes pairs of elementary excitations in the correlated
quasiboson vacuum.  Two electron-hole excitations acting on the reference are
a doubly excited determinant, which turns this pair amplitude into a natural
CC doubles amplitude.  Requiring the direct $\eh$, transverse $\eh$, and
$\pp$ representations to use the same antisymmetric $t_{ij}^{ab}$ then ties
the three channels together.

In this language, the familiar pieces of CCD have a simple origin.  The free
term $\Delta$ comes from the inverse two-particle propagator $L_0^{-1}$ and
appears with the same sign in all three channel representations.  The bare
doubles driver comes from the fully irreducible interaction $\Lambda$.  The
direct $\eh$, transverse $\eh$, and $\pp$ Riccati equations generate the
three interaction contributions $\Xi^r$, which become the ring,
crossed-ring, and ladder terms for bare kernels.  The mosaic terms complete
the one-body part.  Combining these common and channel-specific pieces once
reproduces the projected CCD equation.

The genuine parquet vertices remain the frequency-dependent quantities
$\Phi^r(\omega)$ obtained from the channel BSE resolvents.  They should not
be identified with the Riccati-space quantities $\Xi^r[\bt]$.  A controlled
extension beyond CCD therefore requires either retaining the full dynamical
$\Phi^r$ in the parquet kernel updates or specifying an additional static
projection.  Determining the most useful closure is the next step.

%%%%%%%%%%%%%%%%%%%%%%%%
\bibliography{parquet-mr-hedin}
%%%%%%%%%%%%%%%%%%%%%%%%

\end{document}
